\section{Node/Root Supervision Allocation Simulations}\label{app:v4-supervision-allocation}

The Markov benchmark makes the supervision tradeoff observable with a known
oracle. Each synthetic document has a root transition-count label. Each leaf
or merge span can also receive the same oracle readout on its own span. Those
node labels are sampled with known propensities, matching the audit design in
Section~\ref{sec:min-audit}.

\subsection{Budget Accounting}

The controlled runs compare root-only supervision with reallocated
supervision at a fixed token-observation budget. A full-document label on a
$128$-token sequence costs $128$ token observations. A label on a $16$-token
leaf costs $16$ token observations, so eight such leaf labels have the same
budget mass as one root label. The root-label share ranges from $100\%$ to
$10\%$. Reallocation spends the remaining budget on local labels.

Three local-label policies appear in the figures. \emph{Leaf-only} spends the
local budget on leaf count labels. \emph{Depth-equal} splits local labels
across non-root depths. \emph{Balanced-node} spreads labels across non-root
nodes. Each policy produces randomly assigned node supervision under a known
sampling rule. In a real audit, the same records become logged propensities
for IPW/Horvitz--Thompson estimation.

\subsection{Allocation Patterns}

Figure~\ref{fig:min-markov-leaf-size} reports the compact heatmap view:
tree supervision matches or improves on the flat root-only baseline across
the budget range and remains competitive with only $10\%$ root-label
coverage. Figures~\ref{fig:v4-simple-leaf-mass} and
\ref{fig:v4-hard-leaf-mass} show the same result as per-policy curves. The
recoverable case has a broad safe band of leaf sizes. The structural case is
harder, but the tree still dominates the flat baseline across the shown
budgets.

\begin{figure}[t]
    \centering
    \includegraphics[width=0.95\linewidth]{markov_simple_leaf_mass.png}
    \caption{\textbf{Recoverable Markov supervision allocation.}
    Each panel fixes a root-label share. The green curve uses root labels
    only. The blue, purple, and rose curves keep the $100\%$ token budget and
    reallocate missing root-label mass to local labels under the three
    policies. The amber curve is the flat FNO baseline. Lower MAE is better.}
    \label{fig:v4-simple-leaf-mass}
\end{figure}

\begin{figure}[t]
    \centering
    \includegraphics[width=0.95\linewidth]{markov_hard_leaf_mass.png}
    \caption{\textbf{Structural Markov supervision allocation.}
    The structural regime repeats the same protocol with more registers and
    more expected shifts. Local supervision remains effective, while very
    small leaves expose the boundary-state failure mode.}
    \label{fig:v4-hard-leaf-mass}
\end{figure}

The budget-split view in Figure~\ref{fig:v4-budget-split} isolates the
supervision object shift. The total token budget is held fixed; the horizontal
axis changes the fraction spent on root labels versus local labels. The
reallocation curves stay nearly flat from full root supervision down to about
$40\%$ root supervision before degrading. In this controlled setting,
randomly assigned node labels can replace a large share of root labels because
they measure the same oracle-preservation property.

\begin{figure}[t]
    \centering
    \includegraphics[width=0.95\linewidth]{markov_budget_split.png}
    \caption{\textbf{Splitting a fixed token budget between root and local
    labels.} The green star is the all-root reference. The green line spends
    only the reduced root-label budget. The blue, purple, and rose curves keep
    the total budget fixed by assigning the remaining mass to local labels.
    The flat region shows the range where node supervision preserves root
    accuracy at the same total budget.}
    \label{fig:v4-budget-split}
\end{figure}

\subsection{Interpretation}

The simulations support a narrow methodological claim. Local labels become
substitutes for root labels when they are randomized over the relevant tree
population and measure the same oracle-preservation property. Their value is
statistical as well as operational: they provide training signal for the
compression operators and design-based estimates of local-law distortion. The
Manifesto/Benoit results use the root-observed tier today. Quasi-sentence
codings supply the natural node-level oracle path for the same design.

\clearpage
